\section{可分多项式与可分扩张}\label{sec:galois-2-5}
\subsection{多项式的求导}\label{sec:galois-2-5-1}
设 $A$ 为 $k$-代数. $A$ 的一个\newterm{galois-derivation}{导子}{derivation} $\partial$ 是 $A$ 的一个 $k$-线性自同态，满足公理
\begin{align*}
\forall f,g\in A,\qquad \partial(fg)=\partial(f)g+f\partial(g).
\end{align*}
于是由归纳法可见，对每个 $n\in\mathbb N$ 都有 $\partial(f^n)=n\partial(f)$，特别地，对每个 $\lambda\in k$ 都有 $\partial(\lambda)=\lambda\partial(1)=0$.\footnote{译注：原文的幂公式如此；通常应为 $\partial(f^n)=nf^{n-1}\partial(f)$（$n\geq1$）.}

因此在代数 $A=k[X]$ 上，每个导子都由它在 $X$ 处的值唯一确定. 记 $\partial$ 为 $k[X]$ 上满足 $\partial(X)=1$ 的唯一导子. 对多项式 $f=a_nX^n+a_{n-1}X^{n-1}+\cdots+a_0$，有
\begin{align*}
\partial f=na_nX^{n-1}+(n-1)a_{n-1}X^{n-2}+\cdots+a_1.
\end{align*}
称 $\partial f$ 为 $f$ 的导多项式，无歧义时记为 $f'$.
\begin{remark*}
设 $\tau:k\to k'$ 是域态射. 将延拓 $\tau$ 且把 $X$ 映为 $X$ 的环态射 $k[X]\longrightarrow k'[X]$ 也记为 $\tau$. 则对每个 $f\in k[X]$，都有 $\tau(f)'=\tau(f')$.
\end{remark*}
\begin{lemma*}
设 $f\in k[X]$ 满足 $f'=0$.
\begin{enumerate}[label=\roman*)]
\item 若 $k$ 的特征为零，则 $\deg(f)=0$.
\item 若 $k$ 的特征为 $p>0$，则存在唯一的多项式 $g\in k[X]$ 使得 $f=g(X^p)$.
\end{enumerate}
\end{lemma*}
\begin{proof}
i) 显然. 对于 ii)，写出 $f=\sum_na_nX^n$ 及 $f'=\sum_nna_nX^{n-1}=0$，可见 $a_n\ne0\Rightarrow p\mid n$，所以 $f=\sum_{m\in\mathbb N}a_{pm}X^{pm}=g(X^p)$，其中 $g=\sum_ma_{pm}X^m$. $g$ 的唯一性显然.
\end{proof}
\subsection{可分多项式}\label{sec:galois-2-5-2}
\begin{definition*}
若 $k[X]$ 的理想 $(f,f')$ 是单位理想（即 $f$ 与 $f'$ 互素），则称多项式 $f\in k[X]$ 为\newterm{galois-separable-polynomial}{可分多项式}{separable polynomial}.
\end{definition*}
设 $f\in k[X]$. 若 $\overline k$ 是 $k$ 的代数闭包，则 $f$ 在 $\overline k[X]$ 中分裂为 $f=a_n(X-\alpha_1)^{v_1}\cdots(X-\alpha_m)^{v_m}$，其中 $a_n$ 是 $f$ 的首项系数（即 $n=\deg(f)$），$\sum_{i=1}^m v_i=n$，且假定 $\alpha_i\in\overline k$ 互不相同. 因此 $\alpha_i$ 是 $f$ 在 $\overline k$ 中的根，$v_i$ 是根 $\alpha_i$ 的``重数''.
\begin{proposition*}
对于 $f\in k[X]$，以下性质等价：
\begin{enumerate}[label=\roman*)]
\item $f$ 可分.
\item $f$ 与 $f'$ 在 $k$ 的一个代数闭包中无公共根.
\item $f$ 在 $k$ 的一个代数闭包中的所有根的重数均为 $1$.
\item[iii$'$)] $f$ 在任意代数闭包中都有 $\deg(f)$ 个互异的根.
\item[iii$''$)] 若 $\overline k$ 为 $k$ 的代数闭包，则 $\overline k$-代数 $\overline k[X]/(f)$ 是既约的（此时它同构于 $\overline k^{\deg(f)}=\overline k\times\overline k\times\cdots\times\overline k$）.
\end{enumerate}
\end{proposition*}
\begin{proof}
i)$\Rightarrow$ii). 设 $\overline k$ 为 $k$ 的代数闭包. 若 $g,h\in k[X]$ 满足 $fg+f'h=1$，则在 $\overline k[X]$ 中的同一等式说明 $f$ 与 $f'$ 在其中没有公共不可约因子，因而无公共根.

ii)$\Rightarrow$iii). 证明其逆否命题. 假设 $f$ 在一个代数闭包 $\overline k$ 中有二重根 $\alpha$，则存在 $g\in\overline k[X]$ 使得 $f=(X-\alpha)^2\cdot g$. 从而 $f'=(X-\alpha)(2g+(X-\alpha)g')$，说明 $\alpha$ 为 $f$ 与 $f'$ 的公共根.

iii)$\Rightarrow$i). 仍证明逆否命题. 若 $(f,f')$ 不是 $k[X]$ 的单位理想，则 $f$ 与 $f'$ 有公共不可约因子，记为 $h\in k[X]$. 因此存在 $g\in k[X]$ 使得 $f=hg$. 求导得 $f'=h'g+hg'$. 由于 $h\mid f'$，可知 $h\mid h'g$. 有以下两种可能：
\begin{itemize}
\item 若 $h'\ne0$，则由 $\deg(h')<\deg(h)$ 知 $h$ 不整除 $h'$，故由 Euclid 引理，$h$ 整除 $g$. 因此 $h^2$ 整除 $f$. 而 $h$ 在 $\overline k$ 中有根，此根就是 $f$ 的二重根.
\item 若 $h'=0$，则由上一个引理，$k$ 的特征为 $p>0$，且存在 $e\in k[X]$ 使得 $h=e(X^p)$. 此时 $e$ 在 $\overline k$ 中有根 $\alpha$，故 $X^p-\alpha$ 在 $k[X]$ 中整除 $h$.\footnote{译注：原文此处写 $k[X]$；由 $\alpha\in\overline k$，应在 $\overline k[X]$ 中理解这一整除关系.} 但 $\alpha$ 在 $\overline k$ 中有 $p$ 次根 $\beta$，所以 $X^p-\alpha=(X-\beta)^p$，$\beta$ 为 $h$ 的重根，因而也是 $f$ 的重根.
\end{itemize}
iii) 与 iii$'$) 的等价性显然. 尚需验证 iii)$\Leftrightarrow$iii$''$). 为此，在 $\overline k[X]$ 中将 $f$ 分裂为 $f=a_n(X-\alpha_1)^{v_1}\cdots(X-\alpha_m)^{v_m}$，其中 $\alpha_i$ 两两互异. 由中国剩余定理可知
\begin{align*}
\overline k[X]/(f)=\prod_{i=1}^m\overline k[X]/(X-\alpha_i)^{v_i}.
\end{align*}
这个环既约当且仅当每个因子 $\overline k[X]/(X-\alpha_i)^{v_i}$ 既约，而这等价于 $v_i=1$. 此时 $m=\deg(f)$，故 $\overline k[X]/(f)\simeq\overline k^{\deg(f)}$.
\end{proof}
\begin{remark*}
由性质 iii) 可见，若 $f$ 整除一个可分多项式，则 $f$ 可分.
\end{remark*}
\paragraph{应用（有限域）.}
现在可以完善\hyperref[sec:galois-2-4-4]{有限域定理}. 由于多项式 $X^{p^r}-X$ 可分（其导多项式为 $f'=-1$），它在 $\overline{\mathbb F}_p$ 中有 $p^r$ 个互异的根. 因而由该定理，其分裂域 $\overline{\mathbb F}_p^{F^r}$ 的基数为 $p^r$. 由此得到有限域的分类定理：
\subsection{有限域的分类定理}\label{sec:galois-2-5-3}
\begin{theorem*}
对于任意素数幂 $p^r$，存在基数为 $p^r$ 的域 $\mathbb F_{p^r}$，且在同构意义下唯一. 它是多项式 $X^{p^r}-X$ 在 $\mathbb F_p$ 上的分裂域. 每个有限域都具有这种形式.
\end{theorem*}
\begin{proof}
设 $k$ 为有限域. 它的特征必不为零，记为 $p$. 它在素子域 $\mathbb F_p$ 上的次数有限，记为 $r$，故由\hyperref[sec:galois-2-4-4]{有限域定理}，它是 $X^{p^r}-X$ 的分裂域. 这就给出了唯一性. 如上所述，存在性来自该定理及 $X^{p^r}-X$ 的可分性.
\end{proof}
\subsection{可分扩张}\label{sec:galois-2-5-4}
\begin{definition*}
设 $K\supset k$ 为代数扩张.
\begin{itemize}
\item 若 $\alpha\in K$ 的极小多项式 $f_\alpha\in k[X]$ 可分，则称 $\alpha$ 在 $k$ 上\newterm{galois-separable-element}{可分}{separable element}.
\item 若 $K$ 的每个元素在 $k$ 上可分，则称 $K$ 在 $k$ 上可分，并称此扩张为\newterm{galois-separable-extension}{可分扩张}{separable extension}.
\end{itemize}
\end{definition*}
为给出命题~\ref{sec:galois-2-5-2} 对扩张的对应版本，需要考察：对于扩张，什么概念对应于多项式的根？为此，回顾对不可约的 $f\in k[X]$ 有以下双射：
\begin{align*}
\operatorname{Hom}_{k\text{-alg}}(K[X]/(f),\overline k)\xrightarrow{\sim}\{\alpha\in\overline k,f(\alpha)=0\},
\end{align*}
它由 $\iota\mapsto\iota(\overline X)$ 给出，其中 $\overline X$ 是 $X$ 在 $K[X]/(f)$ 中的像.\footnote{译注：原文两处均写 $K[X]/(f)$，此处应为 $k[X]/(f)$.} 因此，与根的概念相应的是嵌入的概念. 下面的引理说明，正如多项式 $f$ 至多有 $\deg(f)$ 个根，有限扩张 $K\supset k$ 至多有 $[K:k]$ 个嵌入.
\begin{proposition*}
设 $K\supset k$ 为有限扩张，$\overline k$ 为 $k$ 的代数闭包. 则 $K$ 到 $\overline k$ 的 $k$-嵌入数满足
\begin{align*}
|\operatorname{Hom}_{k\text{-alg}}(K,\overline k)|\leqslant[K:k].
\end{align*}
\end{proposition*}
\begin{proof}
标量扩张的泛性质给出双射
\begin{align*}
\operatorname{Hom}_{k\text{-alg}}(K,\overline k)\xrightarrow{\sim}\operatorname{Hom}_{\overline k\text{-alg}}(\overline k\otimes_k K,\overline k),\qquad\iota\longmapsto\tau,
\end{align*}
由等式 $\tau(\lambda\otimes\alpha)=\lambda\iota(\alpha)$ 刻画. 置 $I:=\operatorname{Hom}_{\overline k\text{-alg}}(\overline k\otimes_k K,\overline k)$，考虑乘积态射
\begin{align*}
\Pi\tau:\overline k\otimes_k K\longrightarrow\overline k^I,\qquad\lambda\otimes\alpha\longmapsto(\tau(\lambda\otimes\alpha))_{\tau\in I}.
\end{align*}
下面的引理说明此态射为满射，因此
\begin{align*}
[K:k]=\dim_{\overline k}(\overline k\otimes_k K)\geqslant\dim_{\overline k}(\overline k^I)=|\operatorname{Hom}_{k\text{-alg}}(K,\overline k)|.
\end{align*}
\end{proof}
\begin{lemma*}
设 $A$ 为有限维交换 $\overline k$-代数. 则态射
\begin{align*}
\Pi\tau:A\longrightarrow\overline k^{\operatorname{Hom}_{\overline k\text{-alg}}(A,\overline k)},\qquad a\longmapsto(\tau(a))_{\tau\in\operatorname{Hom}_{\overline k\text{-alg}}(A,\overline k)}
\end{align*}
为满射，且其核为 $A$ 的幂零根.
\end{lemma*}
\begin{proof}
固定一个 $\overline k$-代数态射 $\tau:A\longrightarrow\overline k$. 由于对每个 $\lambda\in\overline k$ 都有 $\tau(\lambda\cdot1_A)=\lambda$，$\tau$ 为满射. 因此 $\mathfrak m:=\operatorname{Ker}(\tau)$ 是 $A$ 的极大理想，且 $\tau$ 分解为 $A\twoheadrightarrow A/\mathfrak m\xrightarrow{\overline\tau}\overline k$，其中 $\overline\tau$ 为双射. 事实上，$\overline\tau$ 由 $\mathfrak m$ 确定. 这是因为复合映射 $\iota_{\mathfrak m}:\overline k\longrightarrow A\longrightarrow A/\mathfrak m$ 使 $A/\mathfrak m$ 成为 $\overline k$ 的有限扩张，故由 $\overline k$ 代数闭知它为同构. 但此时 $\overline\tau\circ\iota_{\mathfrak m}$ 是 $\overline k$-代数 $\overline k$ 的自同构，故为恒等映射. 因此 $\tau$ 与态射 $\tau_{\mathfrak m}:A\twoheadrightarrow A/\mathfrak m\xrightarrow{\iota_{\mathfrak m}^{-1}}\overline k$ 一致. 于是已经证明
\begin{align*}
\operatorname{Hom}_{\overline k\text{-alg}}(A,\overline k)=\{\tau_{\mathfrak m},\mathfrak m\in\operatorname{Max}(A)\},
\end{align*}
且命题中的映射 $\Pi\tau$ 分解为
\begin{align*}
A\longrightarrow\prod_{\mathfrak m\in\operatorname{Max}(A)}A/\mathfrak m\xrightarrow{\sim}\overline k^{\operatorname{Hom}_{\overline k\text{-alg}}(A,\overline k)},
\end{align*}
其中第一箭头为 $a\mapsto(a\bmod\mathfrak m)_{\mathfrak m\in\operatorname{Max}(A)}$，第二箭头为乘积 $\prod_{\mathfrak m}\iota_{\mathfrak m}^{-1}$. 中国剩余定理说明 $\Pi\tau$ 为满射. 此外，它的核为 $\bigcap_{\mathfrak m\in\operatorname{Max}(A)}\mathfrak m$. 由于 $A$ 作为 $A$-模具有有限长度，已知某个有限个极大理想的乘积零化 $A$. 因此 $\bigcap_{\mathfrak m\in\operatorname{Max}(A)}\mathfrak m$ 的某个适当的幂也零化 $A$，即此交为幂零理想，因而包含于 $A$ 的幂零根 $\mathcal N(A)$. 反之，由于商环 $A/\bigcap_{\mathfrak m\in\operatorname{Max}(A)}\mathfrak m$ 既约，也有 $\mathcal N(A)\subset\bigcap_{\mathfrak m\in\operatorname{Max}(A)}\mathfrak m$，从而得等式.
\end{proof}
\begin{theorem*}
设 $K\supset k$ 为有限扩张. 以下条件等价：
\begin{enumerate}[label=\roman*)]
\item $K$ 在 $k$ 上可分.
\item 对 $k$ 的任意代数闭包 $\overline k$ 都有 $|\operatorname{Hom}_{k\text{-alg}}(K,\overline k)|=[K:k]$.
\item 对 $k$ 的任意代数闭包 $\overline k$，$\overline k$-代数 $\overline k\otimes_k K$ 既约（此时它同构于 $\overline k^{[K:k]}=\overline k\times\overline k\times\cdots\times\overline k$）.
\end{enumerate}
\end{theorem*}
\begin{proof}
ii)$\Leftrightarrow$iii) 由上述引理及命题证明中的论证立即得到.

i)$\Rightarrow$ii). 先作如下观察. 设 $K'\subset K$ 为 $K$ 的一个 $k$-子扩张，考虑限制映射
\begin{align*}
\operatorname{Hom}_{k\text{-alg}}(K,\overline k)\longrightarrow\operatorname{Hom}_{k\text{-alg}}(K',\overline k),\qquad\iota\longmapsto\iota|_{K'}.
\end{align*}
命题~\ref{sec:galois-2-2-3} 说明此映射为满射. 此外，它在 $\iota':K'\hookrightarrow\overline k$ 上方的纤维是所有延拓 $\iota'$ 的嵌入 $K\hookrightarrow\overline k$ 所成的集合 $\operatorname{Hom}_{K'\text{-alg},\iota'}(K,\overline k)$，即 $K'$-代数态射的集合，其中 $\overline k$ 配备由 $\iota'$ 给出的 $K'$-代数结构. 因此
\begin{align*}
|\operatorname{Hom}_{k\text{-alg}}(K,\overline k)|=\sum_{\iota'\in\operatorname{Hom}_{k\text{-alg}}(K',\overline k)}|\operatorname{Hom}_{K'\text{-alg},\iota'}(K,\overline k)|.
\end{align*}
特别地，若已知对每个 $\iota$ 都有 $|\operatorname{Hom}_{K'\text{-alg},\iota'}(K,\overline k)|=[K:K']$，且 $|\operatorname{Hom}_{k\text{-alg}}(K',\overline k)|=[K':k]$，则得到
\begin{align*}
|\operatorname{Hom}_{k\text{-alg}}(K,\overline k)|=[K':k][K:K']=[K:k].\tag{*}\label{eq:galois-2-5-4-star}
\end{align*}
这一观察允许对 $K$ 在 $k$ 上的生成元个数 $r$ 作归纳. 若 $r=1$，则 $K=k[\alpha_1]=k[X]/(f_{\alpha_1})$，如上所见，$\operatorname{Hom}_{k\text{-alg}}(K,\overline k)$ 与 $f_{\alpha_1}$ 的根集双射；由 $\alpha_1$ 可分，该集合的基数为 $\deg(f_{\alpha_1})=[K:k]$. 设 $K$ 由 $r$ 个元素 $\alpha_1,\ldots,\alpha_r$ 生成，记 $K':=k[\alpha_1,\ldots,\alpha_{r-1}]$. 由归纳假设，可设 $|\operatorname{Hom}_{k\text{-alg}}(K',\overline k)|=[K':k]$. 此外，由于 $K$ 在 $K'$ 上由（至多）一个元素 $\alpha_r$ 生成，对于每个嵌入 $\iota':K'\hookrightarrow\overline k$ 都有 $|\operatorname{Hom}_{K'\text{-alg},\iota'}(K,\overline k)|=[K:K']$（这样的嵌入使 $\overline k$ 成为 $K'$ 的代数闭包）. 由式~\eqref{eq:galois-2-5-4-star} 得到结论.

iii)$\Rightarrow$i). 在证明这一蕴含之前，注意对每个子扩张 $K'\subset K$，典范的 $\overline k$-代数态射 $\overline k\otimes_k K'\longrightarrow\overline k\otimes_k K$ 都是单射. 事实上，若 $(e_i)_{i\in I}$ 是 $k$-向量空间 $K$ 的一组基，且 $(e_i)_{i\in I'}$ 是 $k$-向量空间 $K'$ 的一组基，则 $(1\otimes e_i)_{i\in I}$ 是 $\overline k$-向量空间 $\overline k\otimes_k K$ 的基，而其子族 $(1\otimes e_i)_{i\in I'}$ 是 $\overline k$-向量空间 $\overline k\otimes_k K'$ 的基. 因此若 $\overline k\otimes_k K$ 既约，则 $\overline k\otimes_k K'$ 也既约. 对任意 $\alpha\in K$，将此应用于 $K'=k[\alpha]$. 则 $\overline k\otimes_k k[\alpha]$ 既约；由于 $k[\alpha]\simeq k[X]/(f_\alpha)$，由命题~\ref{sec:galois-2-5-2} 的 iii$''$) 知 $f_\alpha$ 可分. 由 $\alpha$ 的任意性，$K$ 在 $k$ 上可分.
\end{proof}
\paragraph{应用.}
若 $f\in k[X]$ 可分，则 $k[X]/(f)$ 是 $k$ 的可分扩张. 这由 iii)$\Rightarrow$i) 得到.\footnote{译注：原文此应用未写 $f$ 不可约的假设；一般的可分多项式之商未必是域，故这里的扩张论断需加不可约条件.}
\begin{corollary*}[{由证明得到}]
设 $k\subset K'\subset K$ 为有限扩张塔. 若 $K$ 在 $K'$ 上可分，且 $K'$ 在 $k$ 上可分，则 $K$ 在 $k$ 上可分.
\end{corollary*}
\begin{proof}
重复蕴含 i)$\Rightarrow$ii) 中的论证，得到其证明中的式~\eqref{eq:galois-2-5-4-star}，由此可知 $K$ 在 $k$ 上可分.
\end{proof}
\paragraph{应用.}
由可分元素生成的有限扩张 $K$ 是可分扩张（对生成元个数作归纳）. 特别地，$k[X]$ 中可分多项式的分裂域在 $k$ 上可分.
\begin{remark*}
对无限代数扩张 $K\supset k$，定理中的断言 ii) 没有意义. 但所用论证给出等价关系：
\begin{align*}
K\text{ 在 }k\text{ 上可分}\quad\Longleftrightarrow\quad\overline k\otimes_k K\text{ 既约}.
\end{align*}
\end{remark*}
\subsection{本原元定理}\label{sec:galois-2-5-5}
下面的推论足够引人注目，值得赋予一个生动的名称.
\begin{corollary*}[{本原元定理}]
每个有限可分扩张 $K\supset k$ 都是单生成的（即由一个元素生成）.
\end{corollary*}
\begin{proof}
设 $\alpha\in K$，考虑限制映射
\begin{align*}
\operatorname{Hom}_{k\text{-alg}}(K,\overline k)\longrightarrow\operatorname{Hom}_{k\text{-alg}}(k[\alpha],\overline k),\qquad\iota\longmapsto\iota|_{k[\alpha]}.
\end{align*}
已经看到此映射为满射，定义域的基数为 $[K:k]$（因为假设 $K$ 可分），目标集的基数为 $[k[\alpha]:k]=\deg(f_\alpha)$（因为 $\alpha$ 可分）. 因而有等价关系
\begin{align*}
K=k[\alpha]\ \Longleftrightarrow\ [K:k]=[k[\alpha]:k]\ \Longleftrightarrow\ (\iota\mapsto\iota|_{k[\alpha]}\text{ 为单射}).
\end{align*}
另一方面，映射 $\iota'\mapsto\iota'(\alpha)$ 是从 $\operatorname{Hom}_{k\text{-alg}}(k[\alpha],\overline k)$ 到 $\overline k$ 的单射（事实上，它甚至给出到 $f_\alpha$ 在 $\overline k$ 中的根集的双射）. 所以
\begin{align*}
K=k[\alpha]\ \Longleftrightarrow\ (\iota\mapsto\iota(\alpha)\text{ 为单射})\ \Longleftrightarrow\ \alpha\notin\bigcup_{\iota_1\ne\iota_2}\operatorname{Ker}(\iota_1-\iota_2).
\end{align*}
注意每个 $\operatorname{Ker}(\iota_1-\iota_2)$ 都是 $K$ 的真 $k$-子空间. 当 $k$ 无限时，只需援引下面的引理.

\noindent\textbf{引理.} 设 $k$ 为无限域，$V$ 为有限维 $k$-向量空间. 则有限个真 $k$-子空间之并的补集非空.

\noindent\textbf{引理的证明.} 对 $d=\dim_k(V)$ 作归纳. 当 $d=1$ 时结论显然（而且对有限的 $k$ 也成立）. 设 $d>1$，给定真 $k$-子空间 $V_1,\ldots,V_r$. 可以假设 $V_i$ 为两两不同的超平面. 则对每个 $i>1$，$V_1\cap V_i$ 为 $V_1$ 的真 $k$-子空间，故由归纳假设，存在 $v_1\in V_1\setminus\bigcup_{i>1}V_i$. 选取 $w_1\in V\setminus V_1$ 并考虑仿射直线 $D=w_1+kv_1$. 有 $D\cap V_1=\emptyset$. 此外，对 $i>1$ 有 $|D\cap V_i|<1$.\footnote{译注：原文为严格不等式 $<1$；随后论证实际给出 $\leqslant1$，即至多一个交点.} 事实上，若 $w_1+\lambda v_1$ 和 $w_1+\mu v_1$ 都属于 $V_i$，则 $(\mu-\lambda)v_1\in V_i$，故 $\mu=\lambda$. 因此 $D\cap\bigcup_{i=1}^rV_i$ 为有限集，而直线 $D$ 无限，所以其补集非空. $\square$

尚需处理 $k$ 有限的情形. 为此可设 $k=\mathbb F_p$. 则 $K=\mathbb F_{p^r}$，其中 $r=[K:k]$，而 $K^\times$ 是 $p^r-1$ 次单位根群（即 $X^{p^r}-1$ 的根）.\footnote{译注：括号中的原式为 $X^{p^r}-1$，与前面的 $p^r-1$ 次单位根不符；相应多项式应为 $X^{p^r-1}-1$.} 下面的引理说明此群为循环群. 此群的任意生成元 $\alpha$ 也是 $\mathbb F_{p^r}$ 在 $\mathbb F_p$ 上的生成元.
\end{proof}
\begin{lemma*}
设 $k$ 为域，$G\subset k^\times$ 为 $k^\times$ 的有限子群. 则 $G$ 为循环群.
\end{lemma*}
\begin{proof}
记 $n=|G|$. 要证明 $G$ 中存在阶为 $n$ 的元素. 设 $m$ 为 $G$ 的所有元素的阶的最小公倍数. 有限 Abel 群的结构定理（主理想整环 $\mathbb Z$ 上的挠模）说明，$G$ 中存在阶为 $m$ 的元素 $x$（事实上，可以直接而容易地证明这一点，留作练习）. 因而只需证明 $m=n$. 由定义，对每个 $x\in G$ 都有 $x^m=1$，故 $G$ 由 $m$ 次单位根组成，从而 $|G|=n\leqslant m$. 由于 $m\mid n$，便有 $m=n$.
\end{proof}
\begin{remark*}
不可分的有限扩张未必单生成. 例如取 $K=\mathbb F_p(X,Y)\supset k=\mathbb F_p(X^p,Y^p)$. 容易验证，$X^iY^j$（$0\leqslant i,j<p$）这一族是 $K$ 在 $k$ 上的基，因此 $[K:k]=p^2$. 然而，对每个 $\alpha\in K$ 都有 $\alpha^p\in k$，故 $[k[\alpha]:k]=\deg(f_\alpha)\leqslant p$.
\end{remark*}
\begin{example*}
我们已经看到，域 $\mathbb Q(\sqrt[3]2)$ 在 $\mathbb Q$ 上不正规，但添入 $j$ 后便成为正规扩张（此时得到 $X^2-3$ 的分裂域，它在 $\mathbb Q$ 上的次数为 $6$）.\footnote{译注：原文此处确为 $X^2-3$；相关扩张 $\mathbb Q(j,\sqrt[3]2)$ 实为 $X^3-2$ 的分裂域.} 本原元定理告诉我们此域为单生成，但没有告诉我们如何找到生成元. 后面将介绍寻找生成元的方法，并证明例如 $j+\sqrt[3]2$ 的次数为 $6$，因而生成 $\mathbb Q(j,\sqrt[3]2)$.
\end{example*}



